%File: FLAIRS.tex
\documentclass[letterpaper]{article}
\usepackage{flairs}
\usepackage{times}
\usepackage{helvet}
\usepackage{courier}
\usepackage{graphicx}
\usepackage{subfig}
\usepackage{xcolor}
\frenchspacing
\setlength{\pdfpagewidth}{8.5in}
\setlength{\pdfpageheight}{11in}
\clubpenalty=10000
\widowpenalty=10000

\DeclareMathOperator*{\argmax}{arg\,max}

\pdfinfo{
/Title (Temporal PIG: A Predict-Then-Verify Approach to Resisting the Noisy-TV Problem)
/Author (David Mathews, Larry D. Pyeatt)
/Keywords (reinforcement learning, intrinsic motivation, exploration, noisy-TV problem, information gain)
}

\setcounter{secnumdepth}{0}

\title{Temporal PIG: A Predict-Then-Verify Approach\\to Resisting the Noisy-TV Problem}
\author{David Mathews \and Larry D. Pyeatt\\
Department of Electrical Engineering and Computer Science\\
South Dakota School of Mines and Technology\\
Rapid City, SD 57701\\
\{David.Mathews, Larry.Pyeatt\}@sdsmt.edu
}

\begin{document}
\maketitle

\begin{abstract}
\begin{quote}
Many curiosity-driven exploration methods in reinforcement learning
are fundamentally open-loop: they reward predictions of learning
without verifying that learning occurred.  This makes them vulnerable
to the Noisy-TV problem, where stochastic stimuli trap agents into
endlessly seeking novelty from non-informative distractors.  We
propose a \emph{predict-then-verify} principle: reward only the
measured decrease in model uncertainty after each transition.  We
instantiate this as Temporal PIG (TPIG), a modification of Predicted
Information Gain (PIG) that uses the actual change in KL-divergence
as a temporal-difference reward.  Evaluated on grid-world environments
with varying distractor density, TPIG dramatically outperforms PIG
on distraction rate, exploration speed, and model accuracy,
demonstrating that closed-loop verification of learning is the key to
resisting stochastic distractors.
\end{quote}
\end{abstract}

\section{Introduction}

Reinforcement learning (RL) agents learn by taking actions in an
environment and receiving rewards that guide future
behavior~\cite{Sutton1998,Kaelbling1996RLsurvey}.  When environmental
rewards are sparse or absent, agents must generate their own
\emph{intrinsic} rewards to drive
exploration~\cite{DBLP:journals/corr/abs-1908-06976,FToCFaIM:Jurgen,Oudeyer2007intrinsic}.
Common intrinsic reward methods include count-based
novelty~\cite{DBLP:journals/corr/TangHFSCDSTA16,Bellemare2016counting},
prediction-error
curiosity~\cite{Pathak2017ICM,Rein:VIME}, and information-theoretic
approaches~\cite{10.3389/fncir.2013.00037}.

A fundamental weakness of many intrinsic motivation methods is the
\emph{Noisy-TV problem}~\cite{Burda2019curiosity}.
Consider an environment containing a television with a remote control.
Each button press displays a random, never-before-seen channel.
Novelty-seeking agents become trapped, perpetually pressing the remote
because each press yields a ``new'' state, even though no useful
information about the environment is being learned.

Predicted Information Gain
(PIG)~\cite{10.3389/fncir.2013.00037} selects actions by estimating,
via KL-divergence, how much each action would improve the agent's
transition model.  Although PIG is more principled than raw
novelty-seeking, it still falls prey to the Noisy-TV: the stochastic
transitions from button presses yield high predicted divergence, luring
the agent back to the distractor.

The core issue is that PIG \emph{predicts} information gain but never
\emph{confirms} that learning actually occurred.  We propose a
\textbf{predict-then-verify} approach: compute the predicted
information gain both before and after incorporating a new transition,
and use the \emph{actual decrease} as a TD reward.  We instantiate
this as Temporal PIG (TPIG) and show that it resists the Noisy TV
because stochastic distractors produce large predictions but
negligible measured improvement.

\section{Related Work}

\subsection{Intrinsic Motivation for Exploration}

Intrinsic motivation methods address the sparse-reward problem by
letting agents generate their own learning
signals~\cite{DBLP:journals/corr/abs-1908-06976,FToCFaIM:Jurgen}.
This idea has been realized in several families of algorithms.

\emph{Count-based} methods reward visiting novel states by tracking
visitation counts.  Bellemare et
al.~\cite{Bellemare2016counting} introduced pseudo-counts to unify
count-based exploration with intrinsic motivation in
high-dimensional spaces, and Tang et
al.~\cite{DBLP:journals/corr/TangHFSCDSTA16} extended this approach
using hash functions.  However, count-based methods are directly
vulnerable to the Noisy-TV problem: every channel change registers
as a novel state, producing unbounded reward from a single
distractor.

\emph{Prediction-error} methods use the agent's inability to predict
outcomes as a reward signal.  Pathak et
al.~\cite{Pathak2017ICM} proposed the Intrinsic Curiosity Module
(ICM), which trains a forward dynamics model and uses its prediction
error as intrinsic reward.  Houthooft et
al.~\cite{Rein:VIME} proposed VIME, which maximizes information gain
about environment dynamics using Bayesian neural networks.  Burda et
al.~\cite{Burda2019RND} introduced Random Network Distillation (RND),
which measures novelty as the prediction error of a randomly
initialized fixed network.  While these approaches are more robust
than raw count-based methods, Burda et
al.~\cite{Burda2019curiosity} demonstrated in a large-scale study
that prediction-error agents can still be attracted to inherently
unpredictable transitions such as those produced by a Noisy TV.

\emph{Information-theoretic} methods, including PIG~\cite{10.3389/fncir.2013.00037},
estimate the expected change in the agent's world model using
information-theoretic quantities such as KL-divergence.  Little and
Sommer~\shortcite{10.3389/fncir.2013.00037} showed that PIG-based
exploration learns environment structure faster than reward-free
baselines and transfers well to downstream goal-directed tasks.
M\'{e}nard et al.~\cite{DBLP:journals/corr/abs-2007-13442}
studied active learning for pure exploration in RL, providing
theoretical analysis of sample-efficient exploration strategies.

\subsection{The Noisy-TV Problem}

The Noisy-TV problem was identified as a challenge for
curiosity-driven agents by Burda et
al.~\cite{Burda2019curiosity} and has since become a standard test
of robustness for intrinsic motivation
methods~\cite{DBLP:journals/corr/abs-2008-04388}.  Kovac et
al.~\shortcite{DBLP:journals/corr/abs-2008-04388} proposed
goal-sampling strategies based on learning progress to mitigate
sensitivity to noise.  Ladosz et
al.~\cite{Ladosz_2022} survey the broader exploration literature and
note that most intrinsic motivation methods remain susceptible to
stochastic distractors, identifying this as an open problem.

TPIG can also be viewed as an instantiation of \emph{learning
progress}~\cite{Oudeyer2007intrinsic}, which rewards reductions in
prediction error over time.  TPIG operates directly on PIG's
KL-divergence formulation without auxiliary networks, making it a
minimal modification.  The same predict-then-verify principle could
be applied with other divergence measures such as MSE.

\section{Background}

\subsection{Temporal Difference Learning}

In temporal difference (TD) learning~\cite{Sutton1998}, the agent
maintains Q-values $Q(s,a)$ for each state-action
pair~\cite{Watkins1992Qlearning} and updates them after each
transition $(s,a,s')$ using:
\begin{equation}
  Q(s,a) \leftarrow Q(s,a) + \alpha \left[r + \gamma \max_{a'} Q(s',a') - Q(s,a)\right],
  \label{eq:TDQVF}
\end{equation}
where $r$ is the reward, $\alpha$ is the learning rate, and $\gamma$
is the discount factor~\cite{Sutton1998}.

\subsection{Predicted Information Gain (PIG)}

PIG uses the KL-divergence~\cite{Kullback1951information} to estimate
how much an agent's transition model would change if it took a
particular action.  Let
$\hat{\Theta}_{sa\cdot}$ denote the agent's estimated transition
probabilities from state $s$ under action $a$, and let
$\hat{\Theta}_{sa\cdot}^{s \to s^*}$ denote the updated estimate
after hypothetically observing a transition to $s^*$.  The predicted
information gain is:
\begin{equation}
  \textrm{PIG}(s,a) := \sum_{s^*} \hat{\Theta}_{sas^*}\,
  D_{\textrm{KL}}\!\left(\hat{\Theta}_{sa\cdot}^{s \to s^*}
  \;\|\;\hat{\Theta}_{sa\cdot}\right),
  \label{eq:PIG}
\end{equation}
where $D_{\textrm{KL}}$ is the KL-divergence.  PIG selects the action
$a$ that maximizes this quantity.

\subsection{The Noisy-TV Problem}

When the environment contains a Noisy TV, the stochastic channel
changes produce transitions to many distinct states, each with high
predicted information gain.  PIG's action selection thus repeatedly
favors the TV because the weighted average over all possible channel
outcomes remains large, even though the transitions carry no useful
structural information about the environment.

\section{The Predict-Then-Verify Principle}

TPIG modifies PIG by measuring the \emph{actual} information gained
after each transition rather than relying on predictions alone.  After
observing transition $(s,a,s')$, TPIG computes the predicted
information gain both before and after updating the internal model:
\begin{multline}
  r_{\textrm{TPIG}}(s,a,s') :=
  D_{\textrm{KL}}\!\bigl(\hat{\Theta}_{sa\cdot}^{s \to s',\,t-1}
  \;\|\;\hat{\Theta}_{sa\cdot}^{t-1}\bigr) \\
  - D_{\textrm{KL}}\!\bigl(\hat{\Theta}_{sa\cdot}^{s \to s',\,t}
  \;\|\;\hat{\Theta}_{sa\cdot}^{t}\bigr),
  \label{eq:TPIGReward}
\end{multline}
where superscripts $t{-}1$ and $t$ denote the model before and after
incorporating the observed transition.  The first term is the
predicted information gain \emph{before} the update; the second is the
predicted gain \emph{after}.  If genuine learning occurred, the
remaining information decreases and the reward is positive.

Note that both KL-divergence evaluations reuse the same computation
that PIG already performs, so TPIG adds only a single additional
divergence calculation per time step.
This reward is used to update Q-values via standard TD
learning~\eqref{eq:TDQVF}, with $\alpha = 0.9$ and $\gamma = 0.9$.
Q-values are initialized to 1, providing an optimistic prior that
encourages exploration of all actions before learned rewards dominate.
TPIG selects the action with the highest Q-value:
\begin{equation}
  \textrm{TPIG}(s) = \argmax_a Q(s,a).
  \label{eq:TPIG}
\end{equation}

\begin{figure*}
  \centering
  
\includegraphics[width=0.225\textwidth]{"images/4-TV.pdf"}\hfill
  \includegraphics[width=0.225\textwidth]{"images/1-TV.pdf"}\hfill
  \includegraphics[width=0.225\textwidth]{"images/No-TV.pdf"}
  \caption{Grid-world simulations.  Semi-circles are Noisy TVs;
    arrows indicate one-way walls.  Left: 4-TV ($4\times4$).
    Center: 1-TV ($5\times5$).  Right: No-TV ($6\times6$).}
  \label{Fig:Sim}
\end{figure*}


\subsection{The Closed-Minded Model Problem}

A subtle issue arises in the PIG algorithm when a state-action pair
has been observed only once.  After a single observation of transition
$(s,a,s')$, the agent records the probability of that transition as
100\%.  On subsequent visits to state $s$, PIG predicts that taking
action $a$ will deterministically lead to $s'$, making the predicted
and observed distributions identical.  The KL-divergence returns~0,
and PIG considers the action fully learned, even if the true
transition is stochastic with additional possible outcomes.

This ``closed-minded'' behavior causes PIG to neglect areas it has
visited only once, focusing instead on regions where it has already
observed variance.  Both PIG and TPIG address this by initializing the total observed
transition count for each state-action pair to~1 rather than~0.  This
is equivalent to placing a uniform Dirichlet prior on the transition
probabilities, ensuring that the estimated and predicted distributions
are never identical, so the KL-divergence always indicates some
remaining information to be gained.  As observations accumulate, the
prior is overwhelmed by data and the model converges to the true
distribution.

\subsection{Handling Infinities}

When $\hat{\Theta}_{sa\cdot}$ is empty (no transitions observed),
the KL-divergence returns infinity.  These values are replaced with 0
in the reward calculation.  Negative rewards (which can occur when
discovering a new state temporarily increases uncertainty) are clamped
to 0 to avoid penalizing exploration.  Similarly, any Q-values that
are infinite are normalized to~0 before the TD update is applied.

\subsection{Why TPIG Resists the Noisy TV}

When the agent presses the TV remote, it transitions to a ``new''
channel state.  PIG predicts high information gain because the model
would change substantially.  However, after the update, TPIG
recomputes the predicted gain and finds it is still high because the TV
states are so numerous that each individual observation barely reduces
the remaining uncertainty.  The resulting reward is near zero,
and TPIG's Q-values for TV-related actions decay, redirecting the
agent to parts of the environment where observations actually improve
the model.  More precisely, for a purely stochastic transition with
$k$ equally likely outcomes, the $n$-th observation reduces the
predicted information gain by only $O(1/n)$, so the TPIG reward
vanishes as observations accumulate.  Structured transitions, by
contrast, produce larger reductions that persist until the model
converges.  In summary, PIG takes the action it \emph{predicts} will
teach it the most.  TPIG takes the action \emph{experienced} to yield
the most actual learning, a crucial distinction that makes the
predict-then-verify approach robust to stochastic distractors.

\subsection{Epsilon-Greedy Variants}

We also evaluate $\epsilon$-greedy versions of both algorithms:
EPIG and ETPIG.  With probability $\epsilon$ (set to 10\% unless
otherwise noted), the agent selects a random action instead of the
greedy choice.

\section{Experimental Setup}

\subsection{Grid-World Simulations}

We test on three grid-world environments of increasing size
(\figurename~\ref{Fig:Sim}):
\begin{itemize}
  \item \textbf{4-TV} ($4 \times 4$): Multiple Noisy TVs.
  \item \textbf{1-TV} ($5 \times 5$): A single Noisy TV.
  \item \textbf{No-TV} ($6 \times 6$): No distractors (control).
\end{itemize}

In each environment, actions move the agent in one of four cardinal
directions.  Movement succeeds with 90\% probability; with 10\%
probability the agent remains in place.  Walking into a wall always
leaves the agent stationary.  Activating a Noisy TV increments an unbounded noise variable in the
state, so each activation produces a unique state, simulating a TV
with an inexhaustible supply of channels.  The noise variable does
not affect actual transition dynamics: all TV states share the same
transition probabilities regardless of channel.  The noise is
therefore a pure distractor that inflates the apparent state space
without adding any learnable structure.

All agents start in position~5.  Each episode runs for 1000 time
steps with no goal state.  Every policy is run for 100 independent
trials per simulation.  The source code and simulation data are publicly
available~\cite{TPIG:github}.

\subsection{Metrics}

\begin{itemize}
  \item \textbf{Distraction rate}: Cumulative Noisy-TV activations.
  \item \textbf{Learning rate}: Number of undiscovered state
    transitions remaining.
  \item \textbf{Model accuracy}: KL-divergence between the agent's
    model and the true transition probabilities (excluding
    undiscovered transitions, which would contribute infinity).
\end{itemize}

\subsection{Policies Compared}

Five policies are evaluated: Random Walk (uniform random baseline),
PIG, EPIG ($\epsilon$-greedy PIG), TPIG, and ETPIG ($\epsilon$-greedy
TPIG).

\section{Results}

\subsection{Distraction Rate}

Figures~\ref{Fig:DR4TV} and~\ref{Fig:DR1TV} show
cumulative distractions in the 4-TV and 1-TV environments.  The PIG algorithm is
heavily distracted by the Noisy TVs, with its distraction count
growing steadily throughout the simulation.  The Random Walk baseline,
which selects actions uniformly at random, is distracted the least
simply because it has no bias toward any particular action.

TPIG is initially distracted at a rate comparable to PIG, as it has
not yet accumulated enough experience to distinguish informative
transitions from noise.  However, after a brief learning period, TPIG
sharply reduces its rate of distraction.  In the 1-TV environment
(\figurename~\ref{Fig:DR1TV}), TPIG's distraction rate approaches that
of Random Walk, indicating that it has effectively learned to avoid
the TV.  In the 4-TV environment (\figurename~\ref{Fig:DR4TV}), the
improvement is less dramatic due to the higher density of distractors,
but TPIG still substantially outperforms PIG.
By the final time step, TPIG reduces cumulative distractions by 65\% relative to PIG in the 4-TV environment and by 88\% in the 1-TV environment.

The $\epsilon$-greedy variants EPIG and ETPIG both show reduced
distraction compared to their greedy counterparts, as the random
action selection helps escape distraction loops.  In the No-TV control
environment (not shown), all algorithms produce identical flat lines
at zero since there are no distractors.

\begin{figure}
  \centering
    % GNUPLOT: LaTeX picture with Postscript
  \begingroup
    \makeatletter
    \providecommand\color[2][]{%
      \GenericError{(gnuplot) \space\space\space\@spaces}{%
        Package color not loaded in conjunction with
        terminal option `colourtext'%
      }{See the gnuplot documentation for explanation.%
      }{Either use 'blacktext' in gnuplot or load the package
        color.sty in LaTeX.}%
      \renewcommand\color[2][]{}%
    }%
    \providecommand\includegraphics[2][]{%
      \GenericError{(gnuplot) \space\space\space\@spaces}{%
        Package graphicx or graphics not loaded%
      }{See the gnuplot documentation for explanation.%
      }{The gnuplot epslatex terminal needs graphicx.sty or graphics.sty.}%
      \renewcommand\includegraphics[2][]{}%
    }%
    \providecommand\rotatebox[2]{#2}%
    \@ifundefined{ifGPcolor}{%
      \newif\ifGPcolor
      \GPcolortrue
    }{}%
    \@ifundefined{ifGPblacktext}{%
      \newif\ifGPblacktext
      \GPblacktexttrue
    }{}%
    % define a \g@addto@macro without @ in the name:
    \let\gplgaddtomacro\g@addto@macro
    % define empty templates for all commands taking text:
    \gdef\gplbacktext{}%
    \gdef\gplfronttext{}%
    \makeatother
    \ifGPblacktext
      % no textcolor at all
      \def\colorrgb#1{}%
      \def\colorgray#1{}%
    \else
      % gray or color?
      \ifGPcolor
        \def\colorrgb#1{\color[rgb]{#1}}%
        \def\colorgray#1{\color[gray]{#1}}%
        \expandafter\def\csname LTw\endcsname{\color{white}}%
        \expandafter\def\csname LTb\endcsname{\color{black}}%
        \expandafter\def\csname LTa\endcsname{\color{black}}%
        \expandafter\def\csname LT0\endcsname{\color[rgb]{1,0,0}}%
        \expandafter\def\csname LT1\endcsname{\color[rgb]{0,1,0}}%
        \expandafter\def\csname LT2\endcsname{\color[rgb]{0,0,1}}%
        \expandafter\def\csname LT3\endcsname{\color[rgb]{1,0,1}}%
        \expandafter\def\csname LT4\endcsname{\color[rgb]{0,1,1}}%
        \expandafter\def\csname LT5\endcsname{\color[rgb]{1,1,0}}%
        \expandafter\def\csname LT6\endcsname{\color[rgb]{0,0,0}}%
        \expandafter\def\csname LT7\endcsname{\color[rgb]{1,0.3,0}}%
        \expandafter\def\csname LT8\endcsname{\color[rgb]{0.5,0.5,0.5}}%
      \else
        % gray
        \def\colorrgb#1{\color{black}}%
        \def\colorgray#1{\color[gray]{#1}}%
        \expandafter\def\csname LTw\endcsname{\color{white}}%
        \expandafter\def\csname LTb\endcsname{\color{black}}%
        \expandafter\def\csname LTa\endcsname{\color{black}}%
        \expandafter\def\csname LT0\endcsname{\color{black}}%
        \expandafter\def\csname LT1\endcsname{\color{black}}%
        \expandafter\def\csname LT2\endcsname{\color{black}}%
        \expandafter\def\csname LT3\endcsname{\color{black}}%
        \expandafter\def\csname LT4\endcsname{\color{black}}%
        \expandafter\def\csname LT5\endcsname{\color{black}}%
        \expandafter\def\csname LT6\endcsname{\color{black}}%
        \expandafter\def\csname LT7\endcsname{\color{black}}%
        \expandafter\def\csname LT8\endcsname{\color{black}}%
      \fi
    \fi
      \setlength{\unitlength}{0.0500bp}%
      \ifx\gptboxheight\undefined%
        \newlength{\gptboxheight}%
        \newlength{\gptboxwidth}%
        \newsavebox{\gptboxtext}%
      \fi%
      \setlength{\fboxrule}{0.5pt}%
      \setlength{\fboxsep}{1pt}%
      \definecolor{tbcol}{rgb}{1,1,1}%
  \begin{picture}(4740.00,3600.00)%
      \gplgaddtomacro\gplbacktext{%
        \csname LTb\endcsname%%
        \put(454,343){\makebox(0,0)[r]{\strut{}\scriptsize -50}}%
        \csname LTb\endcsname%%
        \put(454,666){\makebox(0,0)[r]{\strut{}\scriptsize 0}}%
        \csname LTb\endcsname%%
        \put(454,990){\makebox(0,0)[r]{\strut{}\scriptsize 50}}%
        \csname LTb\endcsname%%
        \put(454,1314){\makebox(0,0)[r]{\strut{}\scriptsize 100}}%
        \csname LTb\endcsname%%
        \put(454,1638){\makebox(0,0)[r]{\strut{}\scriptsize 150}}%
        \csname LTb\endcsname%%
        \put(454,1962){\makebox(0,0)[r]{\strut{}\scriptsize 200}}%
        \csname LTb\endcsname%%
        \put(454,2286){\makebox(0,0)[r]{\strut{}\scriptsize 250}}%
        \csname LTb\endcsname%%
        \put(454,2610){\makebox(0,0)[r]{\strut{}\scriptsize 300}}%
        \csname LTb\endcsname%%
        \put(454,2934){\makebox(0,0)[r]{\strut{}\scriptsize 350}}%
        \csname LTb\endcsname%%
        \put(454,3258){\makebox(0,0)[r]{\strut{}\scriptsize 400}}%
        \csname LTb\endcsname%%
        \put(507,235){\makebox(0,0){\strut{}\scriptsize 0}}%
        \csname LTb\endcsname%%
        \put(912,235){\makebox(0,0){\strut{}\scriptsize 100}}%
        \csname LTb\endcsname%%
        \put(1318,235){\makebox(0,0){\strut{}\scriptsize 200}}%
        \csname LTb\endcsname%%
        \put(1723,235){\makebox(0,0){\strut{}\scriptsize 300}}%
        \csname LTb\endcsname%%
        \put(2128,235){\makebox(0,0){\strut{}\scriptsize 400}}%
        \csname LTb\endcsname%%
        \put(2533,235){\makebox(0,0){\strut{}\scriptsize 500}}%
        \csname LTb\endcsname%%
        \put(2938,235){\makebox(0,0){\strut{}\scriptsize 600}}%
        \csname LTb\endcsname%%
        \put(3344,235){\makebox(0,0){\strut{}\scriptsize 700}}%
        \csname LTb\endcsname%%
        \put(3749,235){\makebox(0,0){\strut{}\scriptsize 800}}%
        \csname LTb\endcsname%%
        \put(4154,235){\makebox(0,0){\strut{}\scriptsize 900}}%
        \csname LTb\endcsname%%
        \put(4559,235){\makebox(0,0){\strut{}\scriptsize 1000}}%
      }%
      \gplgaddtomacro\gplfronttext{%
        \csname LTb\endcsname%%
        \put(4196,3153){\makebox(0,0)[r]{\strut{}\scriptsize Random Walk}}%
        \csname LTb\endcsname%%
        \put(4196,3028){\makebox(0,0)[r]{\strut{}\scriptsize PIG}}%
        \csname LTb\endcsname%%
        \put(4196,2903){\makebox(0,0)[r]{\strut{}\scriptsize EPIG}}%
        \csname LTb\endcsname%%
        \put(4196,2778){\makebox(0,0)[r]{\strut{}\scriptsize TPIG}}%
        \csname LTb\endcsname%%
        \put(4196,2653){\makebox(0,0)[r]{\strut{}\scriptsize ETPIG}}%
        \csname LTb\endcsname%%
        \put(93,1800){\rotatebox{-270.00}{\makebox(0,0){\strut{}\footnotesize Cumulative Distractions}}}%
        \csname LTb\endcsname%%
        \put(2533,75){\makebox(0,0){\strut{}\footnotesize TimeStep}}%
        \csname LTb\endcsname%%
        \put(2533,3419){\makebox(0,0){\strut{}\small Distraction Rate: 4-TV Simulation}}%
      }%
      \gplbacktext
      \put(0,0){\includegraphics[width={237.00bp},height={180.00bp}]{gnuplot_figs/Distraction_Rate_4-TV}}%
      \gplfronttext
    \end{picture}%
  \endgroup
  \caption{Distraction rate: 4-TV simulation.  Shaded regions show $\pm 1$ standard deviation over 100 trials.}
  \label{Fig:DR4TV}
\end{figure}

\begin{figure}
  \centering
    % GNUPLOT: LaTeX picture with Postscript
  \begingroup
    \makeatletter
    \providecommand\color[2][]{%
      \GenericError{(gnuplot) \space\space\space\@spaces}{%
        Package color not loaded in conjunction with
        terminal option `colourtext'%
      }{See the gnuplot documentation for explanation.%
      }{Either use 'blacktext' in gnuplot or load the package
        color.sty in LaTeX.}%
      \renewcommand\color[2][]{}%
    }%
    \providecommand\includegraphics[2][]{%
      \GenericError{(gnuplot) \space\space\space\@spaces}{%
        Package graphicx or graphics not loaded%
      }{See the gnuplot documentation for explanation.%
      }{The gnuplot epslatex terminal needs graphicx.sty or graphics.sty.}%
      \renewcommand\includegraphics[2][]{}%
    }%
    \providecommand\rotatebox[2]{#2}%
    \@ifundefined{ifGPcolor}{%
      \newif\ifGPcolor
      \GPcolortrue
    }{}%
    \@ifundefined{ifGPblacktext}{%
      \newif\ifGPblacktext
      \GPblacktexttrue
    }{}%
    % define a \g@addto@macro without @ in the name:
    \let\gplgaddtomacro\g@addto@macro
    % define empty templates for all commands taking text:
    \gdef\gplbacktext{}%
    \gdef\gplfronttext{}%
    \makeatother
    \ifGPblacktext
      % no textcolor at all
      \def\colorrgb#1{}%
      \def\colorgray#1{}%
    \else
      % gray or color?
      \ifGPcolor
        \def\colorrgb#1{\color[rgb]{#1}}%
        \def\colorgray#1{\color[gray]{#1}}%
        \expandafter\def\csname LTw\endcsname{\color{white}}%
        \expandafter\def\csname LTb\endcsname{\color{black}}%
        \expandafter\def\csname LTa\endcsname{\color{black}}%
        \expandafter\def\csname LT0\endcsname{\color[rgb]{1,0,0}}%
        \expandafter\def\csname LT1\endcsname{\color[rgb]{0,1,0}}%
        \expandafter\def\csname LT2\endcsname{\color[rgb]{0,0,1}}%
        \expandafter\def\csname LT3\endcsname{\color[rgb]{1,0,1}}%
        \expandafter\def\csname LT4\endcsname{\color[rgb]{0,1,1}}%
        \expandafter\def\csname LT5\endcsname{\color[rgb]{1,1,0}}%
        \expandafter\def\csname LT6\endcsname{\color[rgb]{0,0,0}}%
        \expandafter\def\csname LT7\endcsname{\color[rgb]{1,0.3,0}}%
        \expandafter\def\csname LT8\endcsname{\color[rgb]{0.5,0.5,0.5}}%
      \else
        % gray
        \def\colorrgb#1{\color{black}}%
        \def\colorgray#1{\color[gray]{#1}}%
        \expandafter\def\csname LTw\endcsname{\color{white}}%
        \expandafter\def\csname LTb\endcsname{\color{black}}%
        \expandafter\def\csname LTa\endcsname{\color{black}}%
        \expandafter\def\csname LT0\endcsname{\color{black}}%
        \expandafter\def\csname LT1\endcsname{\color{black}}%
        \expandafter\def\csname LT2\endcsname{\color{black}}%
        \expandafter\def\csname LT3\endcsname{\color{black}}%
        \expandafter\def\csname LT4\endcsname{\color{black}}%
        \expandafter\def\csname LT5\endcsname{\color{black}}%
        \expandafter\def\csname LT6\endcsname{\color{black}}%
        \expandafter\def\csname LT7\endcsname{\color{black}}%
        \expandafter\def\csname LT8\endcsname{\color{black}}%
      \fi
    \fi
      \setlength{\unitlength}{0.0500bp}%
      \ifx\gptboxheight\undefined%
        \newlength{\gptboxheight}%
        \newlength{\gptboxwidth}%
        \newsavebox{\gptboxtext}%
      \fi%
      \setlength{\fboxrule}{0.5pt}%
      \setlength{\fboxsep}{1pt}%
      \definecolor{tbcol}{rgb}{1,1,1}%
  \begin{picture}(4740.00,3600.00)%
      \gplgaddtomacro\gplbacktext{%
        \csname LTb\endcsname%%
        \put(454,343){\makebox(0,0)[r]{\strut{}\scriptsize -50}}%
        \csname LTb\endcsname%%
        \put(454,828){\makebox(0,0)[r]{\strut{}\scriptsize 0}}%
        \csname LTb\endcsname%%
        \put(454,1314){\makebox(0,0)[r]{\strut{}\scriptsize 50}}%
        \csname LTb\endcsname%%
        \put(454,1800){\makebox(0,0)[r]{\strut{}\scriptsize 100}}%
        \csname LTb\endcsname%%
        \put(454,2286){\makebox(0,0)[r]{\strut{}\scriptsize 150}}%
        \csname LTb\endcsname%%
        \put(454,2772){\makebox(0,0)[r]{\strut{}\scriptsize 200}}%
        \csname LTb\endcsname%%
        \put(454,3258){\makebox(0,0)[r]{\strut{}\scriptsize 250}}%
        \csname LTb\endcsname%%
        \put(507,235){\makebox(0,0){\strut{}\scriptsize 0}}%
        \csname LTb\endcsname%%
        \put(912,235){\makebox(0,0){\strut{}\scriptsize 100}}%
        \csname LTb\endcsname%%
        \put(1318,235){\makebox(0,0){\strut{}\scriptsize 200}}%
        \csname LTb\endcsname%%
        \put(1723,235){\makebox(0,0){\strut{}\scriptsize 300}}%
        \csname LTb\endcsname%%
        \put(2128,235){\makebox(0,0){\strut{}\scriptsize 400}}%
        \csname LTb\endcsname%%
        \put(2533,235){\makebox(0,0){\strut{}\scriptsize 500}}%
        \csname LTb\endcsname%%
        \put(2938,235){\makebox(0,0){\strut{}\scriptsize 600}}%
        \csname LTb\endcsname%%
        \put(3344,235){\makebox(0,0){\strut{}\scriptsize 700}}%
        \csname LTb\endcsname%%
        \put(3749,235){\makebox(0,0){\strut{}\scriptsize 800}}%
        \csname LTb\endcsname%%
        \put(4154,235){\makebox(0,0){\strut{}\scriptsize 900}}%
        \csname LTb\endcsname%%
        \put(4559,235){\makebox(0,0){\strut{}\scriptsize 1000}}%
      }%
      \gplgaddtomacro\gplfronttext{%
        \csname LTb\endcsname%%
        \put(4196,3153){\makebox(0,0)[r]{\strut{}\scriptsize Random Walk}}%
        \csname LTb\endcsname%%
        \put(4196,3028){\makebox(0,0)[r]{\strut{}\scriptsize PIG}}%
        \csname LTb\endcsname%%
        \put(4196,2903){\makebox(0,0)[r]{\strut{}\scriptsize EPIG}}%
        \csname LTb\endcsname%%
        \put(4196,2778){\makebox(0,0)[r]{\strut{}\scriptsize TPIG}}%
        \csname LTb\endcsname%%
        \put(4196,2653){\makebox(0,0)[r]{\strut{}\scriptsize ETPIG}}%
        \csname LTb\endcsname%%
        \put(93,1800){\rotatebox{-270.00}{\makebox(0,0){\strut{}\footnotesize Cumulative Distractions}}}%
        \csname LTb\endcsname%%
        \put(2533,75){\makebox(0,0){\strut{}\footnotesize TimeStep}}%
        \csname LTb\endcsname%%
        \put(2533,3419){\makebox(0,0){\strut{}\small Distraction Rate: 1-TV Simulation}}%
      }%
      \gplbacktext
      \put(0,0){\includegraphics[width={237.00bp},height={180.00bp}]{gnuplot_figs/Distraction_Rate_1-TV}}%
      \gplfronttext
    \end{picture}%
  \endgroup
  \caption{Distraction rate: 1-TV simulation.}
  \label{Fig:DR1TV}
\end{figure}


\subsection{Learning Rate}

Figures~\ref{Fig:MS4TV},~\ref{Fig:MS1TV}, and~\ref{Fig:MS0TV} show
the number of undiscovered state transitions over time.  PIG initially
explores faster than Random Walk in the 1-TV
(\figurename~\ref{Fig:MS1TV}) and No-TV (\figurename~\ref{Fig:MS0TV})
environments, owing to its directed exploration strategy.  However, in
the 4-TV environment (\figurename~\ref{Fig:MS4TV}), PIG's initial
advantage disappears as it becomes trapped by the Noisy TVs,
ultimately performing worse than Random Walk.

TPIG starts slower than both PIG and Random Walk in all three
environments.  This is expected: TPIG must first accumulate enough
transitions to build meaningful Q-values before its directed
exploration becomes effective.  After this initial learning phase,
TPIG consistently surpasses both PIG and Random Walk, discovering
new transitions at a faster rate.  This crossover occurs earlier in
smaller environments and later in larger ones.

The $\epsilon$-greedy variants show no significant impact on learning
rate, indicating that at $\epsilon = 10\%$, the random perturbation
is too small to meaningfully affect the rate of new state discovery.

\begin{figure}
  \centering
    % GNUPLOT: LaTeX picture with Postscript
  \begingroup
    \makeatletter
    \providecommand\color[2][]{%
      \GenericError{(gnuplot) \space\space\space\@spaces}{%
        Package color not loaded in conjunction with
        terminal option `colourtext'%
      }{See the gnuplot documentation for explanation.%
      }{Either use 'blacktext' in gnuplot or load the package
        color.sty in LaTeX.}%
      \renewcommand\color[2][]{}%
    }%
    \providecommand\includegraphics[2][]{%
      \GenericError{(gnuplot) \space\space\space\@spaces}{%
        Package graphicx or graphics not loaded%
      }{See the gnuplot documentation for explanation.%
      }{The gnuplot epslatex terminal needs graphicx.sty or graphics.sty.}%
      \renewcommand\includegraphics[2][]{}%
    }%
    \providecommand\rotatebox[2]{#2}%
    \@ifundefined{ifGPcolor}{%
      \newif\ifGPcolor
      \GPcolortrue
    }{}%
    \@ifundefined{ifGPblacktext}{%
      \newif\ifGPblacktext
      \GPblacktexttrue
    }{}%
    % define a \g@addto@macro without @ in the name:
    \let\gplgaddtomacro\g@addto@macro
    % define empty templates for all commands taking text:
    \gdef\gplbacktext{}%
    \gdef\gplfronttext{}%
    \makeatother
    \ifGPblacktext
      % no textcolor at all
      \def\colorrgb#1{}%
      \def\colorgray#1{}%
    \else
      % gray or color?
      \ifGPcolor
        \def\colorrgb#1{\color[rgb]{#1}}%
        \def\colorgray#1{\color[gray]{#1}}%
        \expandafter\def\csname LTw\endcsname{\color{white}}%
        \expandafter\def\csname LTb\endcsname{\color{black}}%
        \expandafter\def\csname LTa\endcsname{\color{black}}%
        \expandafter\def\csname LT0\endcsname{\color[rgb]{1,0,0}}%
        \expandafter\def\csname LT1\endcsname{\color[rgb]{0,1,0}}%
        \expandafter\def\csname LT2\endcsname{\color[rgb]{0,0,1}}%
        \expandafter\def\csname LT3\endcsname{\color[rgb]{1,0,1}}%
        \expandafter\def\csname LT4\endcsname{\color[rgb]{0,1,1}}%
        \expandafter\def\csname LT5\endcsname{\color[rgb]{1,1,0}}%
        \expandafter\def\csname LT6\endcsname{\color[rgb]{0,0,0}}%
        \expandafter\def\csname LT7\endcsname{\color[rgb]{1,0.3,0}}%
        \expandafter\def\csname LT8\endcsname{\color[rgb]{0.5,0.5,0.5}}%
      \else
        % gray
        \def\colorrgb#1{\color{black}}%
        \def\colorgray#1{\color[gray]{#1}}%
        \expandafter\def\csname LTw\endcsname{\color{white}}%
        \expandafter\def\csname LTb\endcsname{\color{black}}%
        \expandafter\def\csname LTa\endcsname{\color{black}}%
        \expandafter\def\csname LT0\endcsname{\color{black}}%
        \expandafter\def\csname LT1\endcsname{\color{black}}%
        \expandafter\def\csname LT2\endcsname{\color{black}}%
        \expandafter\def\csname LT3\endcsname{\color{black}}%
        \expandafter\def\csname LT4\endcsname{\color{black}}%
        \expandafter\def\csname LT5\endcsname{\color{black}}%
        \expandafter\def\csname LT6\endcsname{\color{black}}%
        \expandafter\def\csname LT7\endcsname{\color{black}}%
        \expandafter\def\csname LT8\endcsname{\color{black}}%
      \fi
    \fi
      \setlength{\unitlength}{0.0500bp}%
      \ifx\gptboxheight\undefined%
        \newlength{\gptboxheight}%
        \newlength{\gptboxwidth}%
        \newsavebox{\gptboxtext}%
      \fi%
      \setlength{\fboxrule}{0.5pt}%
      \setlength{\fboxsep}{1pt}%
      \definecolor{tbcol}{rgb}{1,1,1}%
  \begin{picture}(4740.00,3600.00)%
      \gplgaddtomacro\gplbacktext{%
        \csname LTb\endcsname%%
        \put(454,343){\makebox(0,0)[r]{\strut{}\scriptsize 0}}%
        \csname LTb\endcsname%%
        \put(454,828){\makebox(0,0)[r]{\strut{}\scriptsize 20}}%
        \csname LTb\endcsname%%
        \put(454,1314){\makebox(0,0)[r]{\strut{}\scriptsize 40}}%
        \csname LTb\endcsname%%
        \put(454,1800){\makebox(0,0)[r]{\strut{}\scriptsize 60}}%
        \csname LTb\endcsname%%
        \put(454,2286){\makebox(0,0)[r]{\strut{}\scriptsize 80}}%
        \csname LTb\endcsname%%
        \put(454,2772){\makebox(0,0)[r]{\strut{}\scriptsize 100}}%
        \csname LTb\endcsname%%
        \put(454,3258){\makebox(0,0)[r]{\strut{}\scriptsize 120}}%
        \csname LTb\endcsname%%
        \put(507,235){\makebox(0,0){\strut{}\scriptsize 0}}%
        \csname LTb\endcsname%%
        \put(912,235){\makebox(0,0){\strut{}\scriptsize 100}}%
        \csname LTb\endcsname%%
        \put(1318,235){\makebox(0,0){\strut{}\scriptsize 200}}%
        \csname LTb\endcsname%%
        \put(1723,235){\makebox(0,0){\strut{}\scriptsize 300}}%
        \csname LTb\endcsname%%
        \put(2128,235){\makebox(0,0){\strut{}\scriptsize 400}}%
        \csname LTb\endcsname%%
        \put(2533,235){\makebox(0,0){\strut{}\scriptsize 500}}%
        \csname LTb\endcsname%%
        \put(2938,235){\makebox(0,0){\strut{}\scriptsize 600}}%
        \csname LTb\endcsname%%
        \put(3344,235){\makebox(0,0){\strut{}\scriptsize 700}}%
        \csname LTb\endcsname%%
        \put(3749,235){\makebox(0,0){\strut{}\scriptsize 800}}%
        \csname LTb\endcsname%%
        \put(4154,235){\makebox(0,0){\strut{}\scriptsize 900}}%
        \csname LTb\endcsname%%
        \put(4559,235){\makebox(0,0){\strut{}\scriptsize 1000}}%
      }%
      \gplgaddtomacro\gplfronttext{%
        \csname LTb\endcsname%%
        \put(4196,3153){\makebox(0,0)[r]{\strut{}\scriptsize Random Walk}}%
        \csname LTb\endcsname%%
        \put(4196,3028){\makebox(0,0)[r]{\strut{}\scriptsize PIG}}%
        \csname LTb\endcsname%%
        \put(4196,2903){\makebox(0,0)[r]{\strut{}\scriptsize EPIG}}%
        \csname LTb\endcsname%%
        \put(4196,2778){\makebox(0,0)[r]{\strut{}\scriptsize TPIG}}%
        \csname LTb\endcsname%%
        \put(4196,2653){\makebox(0,0)[r]{\strut{}\scriptsize ETPIG}}%
        \csname LTb\endcsname%%
        \put(93,1800){\rotatebox{-270.00}{\makebox(0,0){\strut{}\footnotesize Undiscovered State Transitions}}}%
        \csname LTb\endcsname%%
        \put(2533,75){\makebox(0,0){\strut{}\footnotesize TimeStep}}%
        \csname LTb\endcsname%%
        \put(2533,3419){\makebox(0,0){\strut{}\small Missed State Transitions: 4-TV Simulation}}%
      }%
      \gplbacktext
      \put(0,0){\includegraphics[width={237.00bp},height={180.00bp}]{gnuplot_figs/Missed_States_4-TV}}%
      \gplfronttext
    \end{picture}%
  \endgroup
  \caption{Missed state transitions: 4-TV simulation.}
  \label{Fig:MS4TV}
\end{figure}

\begin{figure}[t]
  \centering
    % GNUPLOT: LaTeX picture with Postscript
  \begingroup
    \makeatletter
    \providecommand\color[2][]{%
      \GenericError{(gnuplot) \space\space\space\@spaces}{%
        Package color not loaded in conjunction with
        terminal option `colourtext'%
      }{See the gnuplot documentation for explanation.%
      }{Either use 'blacktext' in gnuplot or load the package
        color.sty in LaTeX.}%
      \renewcommand\color[2][]{}%
    }%
    \providecommand\includegraphics[2][]{%
      \GenericError{(gnuplot) \space\space\space\@spaces}{%
        Package graphicx or graphics not loaded%
      }{See the gnuplot documentation for explanation.%
      }{The gnuplot epslatex terminal needs graphicx.sty or graphics.sty.}%
      \renewcommand\includegraphics[2][]{}%
    }%
    \providecommand\rotatebox[2]{#2}%
    \@ifundefined{ifGPcolor}{%
      \newif\ifGPcolor
      \GPcolortrue
    }{}%
    \@ifundefined{ifGPblacktext}{%
      \newif\ifGPblacktext
      \GPblacktexttrue
    }{}%
    % define a \g@addto@macro without @ in the name:
    \let\gplgaddtomacro\g@addto@macro
    % define empty templates for all commands taking text:
    \gdef\gplbacktext{}%
    \gdef\gplfronttext{}%
    \makeatother
    \ifGPblacktext
      % no textcolor at all
      \def\colorrgb#1{}%
      \def\colorgray#1{}%
    \else
      % gray or color?
      \ifGPcolor
        \def\colorrgb#1{\color[rgb]{#1}}%
        \def\colorgray#1{\color[gray]{#1}}%
        \expandafter\def\csname LTw\endcsname{\color{white}}%
        \expandafter\def\csname LTb\endcsname{\color{black}}%
        \expandafter\def\csname LTa\endcsname{\color{black}}%
        \expandafter\def\csname LT0\endcsname{\color[rgb]{1,0,0}}%
        \expandafter\def\csname LT1\endcsname{\color[rgb]{0,1,0}}%
        \expandafter\def\csname LT2\endcsname{\color[rgb]{0,0,1}}%
        \expandafter\def\csname LT3\endcsname{\color[rgb]{1,0,1}}%
        \expandafter\def\csname LT4\endcsname{\color[rgb]{0,1,1}}%
        \expandafter\def\csname LT5\endcsname{\color[rgb]{1,1,0}}%
        \expandafter\def\csname LT6\endcsname{\color[rgb]{0,0,0}}%
        \expandafter\def\csname LT7\endcsname{\color[rgb]{1,0.3,0}}%
        \expandafter\def\csname LT8\endcsname{\color[rgb]{0.5,0.5,0.5}}%
      \else
        % gray
        \def\colorrgb#1{\color{black}}%
        \def\colorgray#1{\color[gray]{#1}}%
        \expandafter\def\csname LTw\endcsname{\color{white}}%
        \expandafter\def\csname LTb\endcsname{\color{black}}%
        \expandafter\def\csname LTa\endcsname{\color{black}}%
        \expandafter\def\csname LT0\endcsname{\color{black}}%
        \expandafter\def\csname LT1\endcsname{\color{black}}%
        \expandafter\def\csname LT2\endcsname{\color{black}}%
        \expandafter\def\csname LT3\endcsname{\color{black}}%
        \expandafter\def\csname LT4\endcsname{\color{black}}%
        \expandafter\def\csname LT5\endcsname{\color{black}}%
        \expandafter\def\csname LT6\endcsname{\color{black}}%
        \expandafter\def\csname LT7\endcsname{\color{black}}%
        \expandafter\def\csname LT8\endcsname{\color{black}}%
      \fi
    \fi
      \setlength{\unitlength}{0.0500bp}%
      \ifx\gptboxheight\undefined%
        \newlength{\gptboxheight}%
        \newlength{\gptboxwidth}%
        \newsavebox{\gptboxtext}%
      \fi%
      \setlength{\fboxrule}{0.5pt}%
      \setlength{\fboxsep}{1pt}%
      \definecolor{tbcol}{rgb}{1,1,1}%
  \begin{picture}(4740.00,3600.00)%
      \gplgaddtomacro\gplbacktext{%
        \csname LTb\endcsname%%
        \put(454,343){\makebox(0,0)[r]{\strut{}\scriptsize 20}}%
        \csname LTb\endcsname%%
        \put(454,707){\makebox(0,0)[r]{\strut{}\scriptsize 40}}%
        \csname LTb\endcsname%%
        \put(454,1071){\makebox(0,0)[r]{\strut{}\scriptsize 60}}%
        \csname LTb\endcsname%%
        \put(454,1436){\makebox(0,0)[r]{\strut{}\scriptsize 80}}%
        \csname LTb\endcsname%%
        \put(454,1800){\makebox(0,0)[r]{\strut{}\scriptsize 100}}%
        \csname LTb\endcsname%%
        \put(454,2165){\makebox(0,0)[r]{\strut{}\scriptsize 120}}%
        \csname LTb\endcsname%%
        \put(454,2529){\makebox(0,0)[r]{\strut{}\scriptsize 140}}%
        \csname LTb\endcsname%%
        \put(454,2893){\makebox(0,0)[r]{\strut{}\scriptsize 160}}%
        \csname LTb\endcsname%%
        \put(454,3258){\makebox(0,0)[r]{\strut{}\scriptsize 180}}%
        \csname LTb\endcsname%%
        \put(507,235){\makebox(0,0){\strut{}\scriptsize 0}}%
        \csname LTb\endcsname%%
        \put(912,235){\makebox(0,0){\strut{}\scriptsize 100}}%
        \csname LTb\endcsname%%
        \put(1318,235){\makebox(0,0){\strut{}\scriptsize 200}}%
        \csname LTb\endcsname%%
        \put(1723,235){\makebox(0,0){\strut{}\scriptsize 300}}%
        \csname LTb\endcsname%%
        \put(2128,235){\makebox(0,0){\strut{}\scriptsize 400}}%
        \csname LTb\endcsname%%
        \put(2533,235){\makebox(0,0){\strut{}\scriptsize 500}}%
        \csname LTb\endcsname%%
        \put(2938,235){\makebox(0,0){\strut{}\scriptsize 600}}%
        \csname LTb\endcsname%%
        \put(3344,235){\makebox(0,0){\strut{}\scriptsize 700}}%
        \csname LTb\endcsname%%
        \put(3749,235){\makebox(0,0){\strut{}\scriptsize 800}}%
        \csname LTb\endcsname%%
        \put(4154,235){\makebox(0,0){\strut{}\scriptsize 900}}%
        \csname LTb\endcsname%%
        \put(4559,235){\makebox(0,0){\strut{}\scriptsize 1000}}%
      }%
      \gplgaddtomacro\gplfronttext{%
        \csname LTb\endcsname%%
        \put(4196,3153){\makebox(0,0)[r]{\strut{}\scriptsize Random Walk}}%
        \csname LTb\endcsname%%
        \put(4196,3028){\makebox(0,0)[r]{\strut{}\scriptsize PIG}}%
        \csname LTb\endcsname%%
        \put(4196,2903){\makebox(0,0)[r]{\strut{}\scriptsize EPIG}}%
        \csname LTb\endcsname%%
        \put(4196,2778){\makebox(0,0)[r]{\strut{}\scriptsize TPIG}}%
        \csname LTb\endcsname%%
        \put(4196,2653){\makebox(0,0)[r]{\strut{}\scriptsize ETPIG}}%
        \csname LTb\endcsname%%
        \put(93,1800){\rotatebox{-270.00}{\makebox(0,0){\strut{}\footnotesize Undiscovered State Transitions}}}%
        \csname LTb\endcsname%%
        \put(2533,75){\makebox(0,0){\strut{}\footnotesize TimeStep}}%
        \csname LTb\endcsname%%
        \put(2533,3419){\makebox(0,0){\strut{}\small Missed State Transitions: 1-TV Simulation}}%
      }%
      \gplbacktext
      \put(0,0){\includegraphics[width={237.00bp},height={180.00bp}]{gnuplot_figs/Missed_States_1-TV}}%
      \gplfronttext
    \end{picture}%
  \endgroup
  \caption{Missed state transitions: 1-TV simulation.}
  \label{Fig:MS1TV}
\end{figure}

\begin{figure}
  \centering
    % GNUPLOT: LaTeX picture with Postscript
  \begingroup
    \makeatletter
    \providecommand\color[2][]{%
      \GenericError{(gnuplot) \space\space\space\@spaces}{%
        Package color not loaded in conjunction with
        terminal option `colourtext'%
      }{See the gnuplot documentation for explanation.%
      }{Either use 'blacktext' in gnuplot or load the package
        color.sty in LaTeX.}%
      \renewcommand\color[2][]{}%
    }%
    \providecommand\includegraphics[2][]{%
      \GenericError{(gnuplot) \space\space\space\@spaces}{%
        Package graphicx or graphics not loaded%
      }{See the gnuplot documentation for explanation.%
      }{The gnuplot epslatex terminal needs graphicx.sty or graphics.sty.}%
      \renewcommand\includegraphics[2][]{}%
    }%
    \providecommand\rotatebox[2]{#2}%
    \@ifundefined{ifGPcolor}{%
      \newif\ifGPcolor
      \GPcolortrue
    }{}%
    \@ifundefined{ifGPblacktext}{%
      \newif\ifGPblacktext
      \GPblacktexttrue
    }{}%
    % define a \g@addto@macro without @ in the name:
    \let\gplgaddtomacro\g@addto@macro
    % define empty templates for all commands taking text:
    \gdef\gplbacktext{}%
    \gdef\gplfronttext{}%
    \makeatother
    \ifGPblacktext
      % no textcolor at all
      \def\colorrgb#1{}%
      \def\colorgray#1{}%
    \else
      % gray or color?
      \ifGPcolor
        \def\colorrgb#1{\color[rgb]{#1}}%
        \def\colorgray#1{\color[gray]{#1}}%
        \expandafter\def\csname LTw\endcsname{\color{white}}%
        \expandafter\def\csname LTb\endcsname{\color{black}}%
        \expandafter\def\csname LTa\endcsname{\color{black}}%
        \expandafter\def\csname LT0\endcsname{\color[rgb]{1,0,0}}%
        \expandafter\def\csname LT1\endcsname{\color[rgb]{0,1,0}}%
        \expandafter\def\csname LT2\endcsname{\color[rgb]{0,0,1}}%
        \expandafter\def\csname LT3\endcsname{\color[rgb]{1,0,1}}%
        \expandafter\def\csname LT4\endcsname{\color[rgb]{0,1,1}}%
        \expandafter\def\csname LT5\endcsname{\color[rgb]{1,1,0}}%
        \expandafter\def\csname LT6\endcsname{\color[rgb]{0,0,0}}%
        \expandafter\def\csname LT7\endcsname{\color[rgb]{1,0.3,0}}%
        \expandafter\def\csname LT8\endcsname{\color[rgb]{0.5,0.5,0.5}}%
      \else
        % gray
        \def\colorrgb#1{\color{black}}%
        \def\colorgray#1{\color[gray]{#1}}%
        \expandafter\def\csname LTw\endcsname{\color{white}}%
        \expandafter\def\csname LTb\endcsname{\color{black}}%
        \expandafter\def\csname LTa\endcsname{\color{black}}%
        \expandafter\def\csname LT0\endcsname{\color{black}}%
        \expandafter\def\csname LT1\endcsname{\color{black}}%
        \expandafter\def\csname LT2\endcsname{\color{black}}%
        \expandafter\def\csname LT3\endcsname{\color{black}}%
        \expandafter\def\csname LT4\endcsname{\color{black}}%
        \expandafter\def\csname LT5\endcsname{\color{black}}%
        \expandafter\def\csname LT6\endcsname{\color{black}}%
        \expandafter\def\csname LT7\endcsname{\color{black}}%
        \expandafter\def\csname LT8\endcsname{\color{black}}%
      \fi
    \fi
      \setlength{\unitlength}{0.0500bp}%
      \ifx\gptboxheight\undefined%
        \newlength{\gptboxheight}%
        \newlength{\gptboxwidth}%
        \newsavebox{\gptboxtext}%
      \fi%
      \setlength{\fboxrule}{0.5pt}%
      \setlength{\fboxsep}{1pt}%
      \definecolor{tbcol}{rgb}{1,1,1}%
  \begin{picture}(4740.00,3600.00)%
      \gplgaddtomacro\gplbacktext{%
        \csname LTb\endcsname%%
        \put(454,343){\makebox(0,0)[r]{\strut{}\scriptsize 50}}%
        \csname LTb\endcsname%%
        \put(454,926){\makebox(0,0)[r]{\strut{}\scriptsize 100}}%
        \csname LTb\endcsname%%
        \put(454,1509){\makebox(0,0)[r]{\strut{}\scriptsize 150}}%
        \csname LTb\endcsname%%
        \put(454,2092){\makebox(0,0)[r]{\strut{}\scriptsize 200}}%
        \csname LTb\endcsname%%
        \put(454,2675){\makebox(0,0)[r]{\strut{}\scriptsize 250}}%
        \csname LTb\endcsname%%
        \put(454,3258){\makebox(0,0)[r]{\strut{}\scriptsize 300}}%
        \csname LTb\endcsname%%
        \put(507,235){\makebox(0,0){\strut{}\scriptsize 0}}%
        \csname LTb\endcsname%%
        \put(912,235){\makebox(0,0){\strut{}\scriptsize 100}}%
        \csname LTb\endcsname%%
        \put(1318,235){\makebox(0,0){\strut{}\scriptsize 200}}%
        \csname LTb\endcsname%%
        \put(1723,235){\makebox(0,0){\strut{}\scriptsize 300}}%
        \csname LTb\endcsname%%
        \put(2128,235){\makebox(0,0){\strut{}\scriptsize 400}}%
        \csname LTb\endcsname%%
        \put(2533,235){\makebox(0,0){\strut{}\scriptsize 500}}%
        \csname LTb\endcsname%%
        \put(2938,235){\makebox(0,0){\strut{}\scriptsize 600}}%
        \csname LTb\endcsname%%
        \put(3344,235){\makebox(0,0){\strut{}\scriptsize 700}}%
        \csname LTb\endcsname%%
        \put(3749,235){\makebox(0,0){\strut{}\scriptsize 800}}%
        \csname LTb\endcsname%%
        \put(4154,235){\makebox(0,0){\strut{}\scriptsize 900}}%
        \csname LTb\endcsname%%
        \put(4559,235){\makebox(0,0){\strut{}\scriptsize 1000}}%
      }%
      \gplgaddtomacro\gplfronttext{%
        \csname LTb\endcsname%%
        \put(4196,3153){\makebox(0,0)[r]{\strut{}\scriptsize Random Walk}}%
        \csname LTb\endcsname%%
        \put(4196,3028){\makebox(0,0)[r]{\strut{}\scriptsize PIG}}%
        \csname LTb\endcsname%%
        \put(4196,2903){\makebox(0,0)[r]{\strut{}\scriptsize EPIG}}%
        \csname LTb\endcsname%%
        \put(4196,2778){\makebox(0,0)[r]{\strut{}\scriptsize TPIG}}%
        \csname LTb\endcsname%%
        \put(4196,2653){\makebox(0,0)[r]{\strut{}\scriptsize ETPIG}}%
        \csname LTb\endcsname%%
        \put(93,1800){\rotatebox{-270.00}{\makebox(0,0){\strut{}\footnotesize Undiscovered State Transitions}}}%
        \csname LTb\endcsname%%
        \put(2533,75){\makebox(0,0){\strut{}\footnotesize TimeStep}}%
        \csname LTb\endcsname%%
        \put(2533,3419){\makebox(0,0){\strut{}\small Missed State Transitions: No-TV Simulation}}%
      }%
      \gplbacktext
      \put(0,0){\includegraphics[width={237.00bp},height={180.00bp}]{gnuplot_figs/Missed_States_No-TV}}%
      \gplfronttext
    \end{picture}%
  \endgroup
  \caption{Missed state transitions: No-TV simulation.}
  \label{Fig:MS0TV}
\end{figure}

\subsection{Model Accuracy}

Figures~\ref{Fig:MA4TV},~\ref{Fig:MA1TV}, and~\ref{Fig:MA0TV} show
the KL-divergence between each agent's estimated model and the true
transition model.  Interpreting this metric requires care: because
undiscovered transitions would contribute infinity to the
KL-divergence, they are excluded from the summation.  Consequently,
an agent that has discovered more states may report a higher
KL-divergence simply because more terms are included.  Model accuracy
must therefore be evaluated jointly with the learning rate results.

In the 4-TV environment (\figurename~\ref{Fig:MA4TV}), PIG shows
worse KL-divergence than Random Walk and has discovered fewer states
by the end of the simulation, placing it clearly behind the baseline.
In the 1-TV environment (\figurename~\ref{Fig:MA1TV}), the comparison
is closer: PIG has a slightly worse KL-divergence but has discovered
more states, making the two roughly equivalent.

TPIG achieves lower KL-divergence than both PIG and Random Walk at
all time steps across all three environments, while simultaneously
having discovered more transitions.  This combination indicates that
TPIG both explores faster and learns the transition structure more
accurately.  The advantage is especially pronounced in the 4-TV
environment, where PIG wastes substantial time on non-informative TV
interactions.

In the No-TV control (\figurename~\ref{Fig:MA0TV}), where there are
no distractors, PIG outperforms Random Walk and TPIG performs best.
This confirms that TPIG's advantage is not merely an artifact of
avoiding distractors; it produces genuinely better exploration even
in clean environments.  The $\epsilon$-greedy variants perform
comparably to their base algorithms, except in the No-TV case where
random perturbation slightly hurts performance by diverting the agent
from its optimal exploration path.
Across all three metrics, TPIG's closed-loop verification
consistently outperforms PIG's open-loop predictions in both noisy
and clean environments.



\begin{figure}[t]
  \centering
    % GNUPLOT: LaTeX picture with Postscript
  \begingroup
    \makeatletter
    \providecommand\color[2][]{%
      \GenericError{(gnuplot) \space\space\space\@spaces}{%
        Package color not loaded in conjunction with
        terminal option `colourtext'%
      }{See the gnuplot documentation for explanation.%
      }{Either use 'blacktext' in gnuplot or load the package
        color.sty in LaTeX.}%
      \renewcommand\color[2][]{}%
    }%
    \providecommand\includegraphics[2][]{%
      \GenericError{(gnuplot) \space\space\space\@spaces}{%
        Package graphicx or graphics not loaded%
      }{See the gnuplot documentation for explanation.%
      }{The gnuplot epslatex terminal needs graphicx.sty or graphics.sty.}%
      \renewcommand\includegraphics[2][]{}%
    }%
    \providecommand\rotatebox[2]{#2}%
    \@ifundefined{ifGPcolor}{%
      \newif\ifGPcolor
      \GPcolortrue
    }{}%
    \@ifundefined{ifGPblacktext}{%
      \newif\ifGPblacktext
      \GPblacktexttrue
    }{}%
    % define a \g@addto@macro without @ in the name:
    \let\gplgaddtomacro\g@addto@macro
    % define empty templates for all commands taking text:
    \gdef\gplbacktext{}%
    \gdef\gplfronttext{}%
    \makeatother
    \ifGPblacktext
      % no textcolor at all
      \def\colorrgb#1{}%
      \def\colorgray#1{}%
    \else
      % gray or color?
      \ifGPcolor
        \def\colorrgb#1{\color[rgb]{#1}}%
        \def\colorgray#1{\color[gray]{#1}}%
        \expandafter\def\csname LTw\endcsname{\color{white}}%
        \expandafter\def\csname LTb\endcsname{\color{black}}%
        \expandafter\def\csname LTa\endcsname{\color{black}}%
        \expandafter\def\csname LT0\endcsname{\color[rgb]{1,0,0}}%
        \expandafter\def\csname LT1\endcsname{\color[rgb]{0,1,0}}%
        \expandafter\def\csname LT2\endcsname{\color[rgb]{0,0,1}}%
        \expandafter\def\csname LT3\endcsname{\color[rgb]{1,0,1}}%
        \expandafter\def\csname LT4\endcsname{\color[rgb]{0,1,1}}%
        \expandafter\def\csname LT5\endcsname{\color[rgb]{1,1,0}}%
        \expandafter\def\csname LT6\endcsname{\color[rgb]{0,0,0}}%
        \expandafter\def\csname LT7\endcsname{\color[rgb]{1,0.3,0}}%
        \expandafter\def\csname LT8\endcsname{\color[rgb]{0.5,0.5,0.5}}%
      \else
        % gray
        \def\colorrgb#1{\color{black}}%
        \def\colorgray#1{\color[gray]{#1}}%
        \expandafter\def\csname LTw\endcsname{\color{white}}%
        \expandafter\def\csname LTb\endcsname{\color{black}}%
        \expandafter\def\csname LTa\endcsname{\color{black}}%
        \expandafter\def\csname LT0\endcsname{\color{black}}%
        \expandafter\def\csname LT1\endcsname{\color{black}}%
        \expandafter\def\csname LT2\endcsname{\color{black}}%
        \expandafter\def\csname LT3\endcsname{\color{black}}%
        \expandafter\def\csname LT4\endcsname{\color{black}}%
        \expandafter\def\csname LT5\endcsname{\color{black}}%
        \expandafter\def\csname LT6\endcsname{\color{black}}%
        \expandafter\def\csname LT7\endcsname{\color{black}}%
        \expandafter\def\csname LT8\endcsname{\color{black}}%
      \fi
    \fi
      \setlength{\unitlength}{0.0500bp}%
      \ifx\gptboxheight\undefined%
        \newlength{\gptboxheight}%
        \newlength{\gptboxwidth}%
        \newsavebox{\gptboxtext}%
      \fi%
      \setlength{\fboxrule}{0.5pt}%
      \setlength{\fboxsep}{1pt}%
      \definecolor{tbcol}{rgb}{1,1,1}%
  \begin{picture}(4740.00,3600.00)%
      \gplgaddtomacro\gplbacktext{%
        \csname LTb\endcsname%%
        \put(400,343){\makebox(0,0)[r]{\strut{}\scriptsize 0}}%
        \csname LTb\endcsname%%
        \put(400,666){\makebox(0,0)[r]{\strut{}\scriptsize 5}}%
        \csname LTb\endcsname%%
        \put(400,990){\makebox(0,0)[r]{\strut{}\scriptsize 10}}%
        \csname LTb\endcsname%%
        \put(400,1314){\makebox(0,0)[r]{\strut{}\scriptsize 15}}%
        \csname LTb\endcsname%%
        \put(400,1638){\makebox(0,0)[r]{\strut{}\scriptsize 20}}%
        \csname LTb\endcsname%%
        \put(400,1962){\makebox(0,0)[r]{\strut{}\scriptsize 25}}%
        \csname LTb\endcsname%%
        \put(400,2286){\makebox(0,0)[r]{\strut{}\scriptsize 30}}%
        \csname LTb\endcsname%%
        \put(400,2610){\makebox(0,0)[r]{\strut{}\scriptsize 35}}%
        \csname LTb\endcsname%%
        \put(400,2934){\makebox(0,0)[r]{\strut{}\scriptsize 40}}%
        \csname LTb\endcsname%%
        \put(400,3258){\makebox(0,0)[r]{\strut{}\scriptsize 45}}%
        \csname LTb\endcsname%%
        \put(454,235){\makebox(0,0){\strut{}\scriptsize 0}}%
        \csname LTb\endcsname%%
        \put(864,235){\makebox(0,0){\strut{}\scriptsize 100}}%
        \csname LTb\endcsname%%
        \put(1275,235){\makebox(0,0){\strut{}\scriptsize 200}}%
        \csname LTb\endcsname%%
        \put(1685,235){\makebox(0,0){\strut{}\scriptsize 300}}%
        \csname LTb\endcsname%%
        \put(2096,235){\makebox(0,0){\strut{}\scriptsize 400}}%
        \csname LTb\endcsname%%
        \put(2507,235){\makebox(0,0){\strut{}\scriptsize 500}}%
        \csname LTb\endcsname%%
        \put(2917,235){\makebox(0,0){\strut{}\scriptsize 600}}%
        \csname LTb\endcsname%%
        \put(3328,235){\makebox(0,0){\strut{}\scriptsize 700}}%
        \csname LTb\endcsname%%
        \put(3738,235){\makebox(0,0){\strut{}\scriptsize 800}}%
        \csname LTb\endcsname%%
        \put(4149,235){\makebox(0,0){\strut{}\scriptsize 900}}%
        \csname LTb\endcsname%%
        \put(4559,235){\makebox(0,0){\strut{}\scriptsize 1000}}%
      }%
      \gplgaddtomacro\gplfronttext{%
        \csname LTb\endcsname%%
        \put(4196,3153){\makebox(0,0)[r]{\strut{}\scriptsize Random Walk}}%
        \csname LTb\endcsname%%
        \put(4196,3028){\makebox(0,0)[r]{\strut{}\scriptsize PIG}}%
        \csname LTb\endcsname%%
        \put(4196,2903){\makebox(0,0)[r]{\strut{}\scriptsize EPIG}}%
        \csname LTb\endcsname%%
        \put(4196,2778){\makebox(0,0)[r]{\strut{}\scriptsize TPIG}}%
        \csname LTb\endcsname%%
        \put(4196,2653){\makebox(0,0)[r]{\strut{}\scriptsize ETPIG}}%
        \csname LTb\endcsname%%
        \put(93,1800){\rotatebox{-270.00}{\makebox(0,0){\strut{}\footnotesize KL-Divergence}}}%
        \csname LTb\endcsname%%
        \put(2507,75){\makebox(0,0){\strut{}\footnotesize TimeStep}}%
        \csname LTb\endcsname%%
        \put(2507,3419){\makebox(0,0){\strut{}\small Model Accuracy: 4-TV Simulation}}%
      }%
      \gplbacktext
      \put(0,0){\includegraphics[width={237.00bp},height={180.00bp}]{gnuplot_figs/Model_Accuracy_4-TV}}%
      \gplfronttext
    \end{picture}%
  \endgroup
  \caption{Model accuracy (KL-divergence): 4-TV simulation.}
  \label{Fig:MA4TV}
\end{figure}

\begin{figure}
  \centering
    % GNUPLOT: LaTeX picture with Postscript
  \begingroup
    \makeatletter
    \providecommand\color[2][]{%
      \GenericError{(gnuplot) \space\space\space\@spaces}{%
        Package color not loaded in conjunction with
        terminal option `colourtext'%
      }{See the gnuplot documentation for explanation.%
      }{Either use 'blacktext' in gnuplot or load the package
        color.sty in LaTeX.}%
      \renewcommand\color[2][]{}%
    }%
    \providecommand\includegraphics[2][]{%
      \GenericError{(gnuplot) \space\space\space\@spaces}{%
        Package graphicx or graphics not loaded%
      }{See the gnuplot documentation for explanation.%
      }{The gnuplot epslatex terminal needs graphicx.sty or graphics.sty.}%
      \renewcommand\includegraphics[2][]{}%
    }%
    \providecommand\rotatebox[2]{#2}%
    \@ifundefined{ifGPcolor}{%
      \newif\ifGPcolor
      \GPcolortrue
    }{}%
    \@ifundefined{ifGPblacktext}{%
      \newif\ifGPblacktext
      \GPblacktexttrue
    }{}%
    % define a \g@addto@macro without @ in the name:
    \let\gplgaddtomacro\g@addto@macro
    % define empty templates for all commands taking text:
    \gdef\gplbacktext{}%
    \gdef\gplfronttext{}%
    \makeatother
    \ifGPblacktext
      % no textcolor at all
      \def\colorrgb#1{}%
      \def\colorgray#1{}%
    \else
      % gray or color?
      \ifGPcolor
        \def\colorrgb#1{\color[rgb]{#1}}%
        \def\colorgray#1{\color[gray]{#1}}%
        \expandafter\def\csname LTw\endcsname{\color{white}}%
        \expandafter\def\csname LTb\endcsname{\color{black}}%
        \expandafter\def\csname LTa\endcsname{\color{black}}%
        \expandafter\def\csname LT0\endcsname{\color[rgb]{1,0,0}}%
        \expandafter\def\csname LT1\endcsname{\color[rgb]{0,1,0}}%
        \expandafter\def\csname LT2\endcsname{\color[rgb]{0,0,1}}%
        \expandafter\def\csname LT3\endcsname{\color[rgb]{1,0,1}}%
        \expandafter\def\csname LT4\endcsname{\color[rgb]{0,1,1}}%
        \expandafter\def\csname LT5\endcsname{\color[rgb]{1,1,0}}%
        \expandafter\def\csname LT6\endcsname{\color[rgb]{0,0,0}}%
        \expandafter\def\csname LT7\endcsname{\color[rgb]{1,0.3,0}}%
        \expandafter\def\csname LT8\endcsname{\color[rgb]{0.5,0.5,0.5}}%
      \else
        % gray
        \def\colorrgb#1{\color{black}}%
        \def\colorgray#1{\color[gray]{#1}}%
        \expandafter\def\csname LTw\endcsname{\color{white}}%
        \expandafter\def\csname LTb\endcsname{\color{black}}%
        \expandafter\def\csname LTa\endcsname{\color{black}}%
        \expandafter\def\csname LT0\endcsname{\color{black}}%
        \expandafter\def\csname LT1\endcsname{\color{black}}%
        \expandafter\def\csname LT2\endcsname{\color{black}}%
        \expandafter\def\csname LT3\endcsname{\color{black}}%
        \expandafter\def\csname LT4\endcsname{\color{black}}%
        \expandafter\def\csname LT5\endcsname{\color{black}}%
        \expandafter\def\csname LT6\endcsname{\color{black}}%
        \expandafter\def\csname LT7\endcsname{\color{black}}%
        \expandafter\def\csname LT8\endcsname{\color{black}}%
      \fi
    \fi
      \setlength{\unitlength}{0.0500bp}%
      \ifx\gptboxheight\undefined%
        \newlength{\gptboxheight}%
        \newlength{\gptboxwidth}%
        \newsavebox{\gptboxtext}%
      \fi%
      \setlength{\fboxrule}{0.5pt}%
      \setlength{\fboxsep}{1pt}%
      \definecolor{tbcol}{rgb}{1,1,1}%
  \begin{picture}(4740.00,3600.00)%
      \gplgaddtomacro\gplbacktext{%
        \csname LTb\endcsname%%
        \put(400,343){\makebox(0,0)[r]{\strut{}\scriptsize 0}}%
        \csname LTb\endcsname%%
        \put(400,828){\makebox(0,0)[r]{\strut{}\scriptsize 10}}%
        \csname LTb\endcsname%%
        \put(400,1314){\makebox(0,0)[r]{\strut{}\scriptsize 20}}%
        \csname LTb\endcsname%%
        \put(400,1800){\makebox(0,0)[r]{\strut{}\scriptsize 30}}%
        \csname LTb\endcsname%%
        \put(400,2286){\makebox(0,0)[r]{\strut{}\scriptsize 40}}%
        \csname LTb\endcsname%%
        \put(400,2772){\makebox(0,0)[r]{\strut{}\scriptsize 50}}%
        \csname LTb\endcsname%%
        \put(400,3258){\makebox(0,0)[r]{\strut{}\scriptsize 60}}%
        \csname LTb\endcsname%%
        \put(454,235){\makebox(0,0){\strut{}\scriptsize 0}}%
        \csname LTb\endcsname%%
        \put(864,235){\makebox(0,0){\strut{}\scriptsize 100}}%
        \csname LTb\endcsname%%
        \put(1275,235){\makebox(0,0){\strut{}\scriptsize 200}}%
        \csname LTb\endcsname%%
        \put(1685,235){\makebox(0,0){\strut{}\scriptsize 300}}%
        \csname LTb\endcsname%%
        \put(2096,235){\makebox(0,0){\strut{}\scriptsize 400}}%
        \csname LTb\endcsname%%
        \put(2507,235){\makebox(0,0){\strut{}\scriptsize 500}}%
        \csname LTb\endcsname%%
        \put(2917,235){\makebox(0,0){\strut{}\scriptsize 600}}%
        \csname LTb\endcsname%%
        \put(3328,235){\makebox(0,0){\strut{}\scriptsize 700}}%
        \csname LTb\endcsname%%
        \put(3738,235){\makebox(0,0){\strut{}\scriptsize 800}}%
        \csname LTb\endcsname%%
        \put(4149,235){\makebox(0,0){\strut{}\scriptsize 900}}%
        \csname LTb\endcsname%%
        \put(4559,235){\makebox(0,0){\strut{}\scriptsize 1000}}%
      }%
      \gplgaddtomacro\gplfronttext{%
        \csname LTb\endcsname%%
        \put(4196,3153){\makebox(0,0)[r]{\strut{}\scriptsize Random Walk}}%
        \csname LTb\endcsname%%
        \put(4196,3028){\makebox(0,0)[r]{\strut{}\scriptsize PIG}}%
        \csname LTb\endcsname%%
        \put(4196,2903){\makebox(0,0)[r]{\strut{}\scriptsize EPIG}}%
        \csname LTb\endcsname%%
        \put(4196,2778){\makebox(0,0)[r]{\strut{}\scriptsize TPIG}}%
        \csname LTb\endcsname%%
        \put(4196,2653){\makebox(0,0)[r]{\strut{}\scriptsize ETPIG}}%
        \csname LTb\endcsname%%
        \put(93,1800){\rotatebox{-270.00}{\makebox(0,0){\strut{}\footnotesize KL-Divergence}}}%
        \csname LTb\endcsname%%
        \put(2507,75){\makebox(0,0){\strut{}\footnotesize TimeStep}}%
        \csname LTb\endcsname%%
        \put(2507,3419){\makebox(0,0){\strut{}\small Model Accuracy: 1-TV Simulation}}%
      }%
      \gplbacktext
      \put(0,0){\includegraphics[width={237.00bp},height={180.00bp}]{gnuplot_figs/Model_Accuracy_1-TV}}%
      \gplfronttext
    \end{picture}%
  \endgroup
  \caption{Model accuracy (KL-divergence): 1-TV simulation.}
  \label{Fig:MA1TV}
\end{figure}

\begin{figure}
  \centering
    % GNUPLOT: LaTeX picture with Postscript
  \begingroup
    \makeatletter
    \providecommand\color[2][]{%
      \GenericError{(gnuplot) \space\space\space\@spaces}{%
        Package color not loaded in conjunction with
        terminal option `colourtext'%
      }{See the gnuplot documentation for explanation.%
      }{Either use 'blacktext' in gnuplot or load the package
        color.sty in LaTeX.}%
      \renewcommand\color[2][]{}%
    }%
    \providecommand\includegraphics[2][]{%
      \GenericError{(gnuplot) \space\space\space\@spaces}{%
        Package graphicx or graphics not loaded%
      }{See the gnuplot documentation for explanation.%
      }{The gnuplot epslatex terminal needs graphicx.sty or graphics.sty.}%
      \renewcommand\includegraphics[2][]{}%
    }%
    \providecommand\rotatebox[2]{#2}%
    \@ifundefined{ifGPcolor}{%
      \newif\ifGPcolor
      \GPcolortrue
    }{}%
    \@ifundefined{ifGPblacktext}{%
      \newif\ifGPblacktext
      \GPblacktexttrue
    }{}%
    % define a \g@addto@macro without @ in the name:
    \let\gplgaddtomacro\g@addto@macro
    % define empty templates for all commands taking text:
    \gdef\gplbacktext{}%
    \gdef\gplfronttext{}%
    \makeatother
    \ifGPblacktext
      % no textcolor at all
      \def\colorrgb#1{}%
      \def\colorgray#1{}%
    \else
      % gray or color?
      \ifGPcolor
        \def\colorrgb#1{\color[rgb]{#1}}%
        \def\colorgray#1{\color[gray]{#1}}%
        \expandafter\def\csname LTw\endcsname{\color{white}}%
        \expandafter\def\csname LTb\endcsname{\color{black}}%
        \expandafter\def\csname LTa\endcsname{\color{black}}%
        \expandafter\def\csname LT0\endcsname{\color[rgb]{1,0,0}}%
        \expandafter\def\csname LT1\endcsname{\color[rgb]{0,1,0}}%
        \expandafter\def\csname LT2\endcsname{\color[rgb]{0,0,1}}%
        \expandafter\def\csname LT3\endcsname{\color[rgb]{1,0,1}}%
        \expandafter\def\csname LT4\endcsname{\color[rgb]{0,1,1}}%
        \expandafter\def\csname LT5\endcsname{\color[rgb]{1,1,0}}%
        \expandafter\def\csname LT6\endcsname{\color[rgb]{0,0,0}}%
        \expandafter\def\csname LT7\endcsname{\color[rgb]{1,0.3,0}}%
        \expandafter\def\csname LT8\endcsname{\color[rgb]{0.5,0.5,0.5}}%
      \else
        % gray
        \def\colorrgb#1{\color{black}}%
        \def\colorgray#1{\color[gray]{#1}}%
        \expandafter\def\csname LTw\endcsname{\color{white}}%
        \expandafter\def\csname LTb\endcsname{\color{black}}%
        \expandafter\def\csname LTa\endcsname{\color{black}}%
        \expandafter\def\csname LT0\endcsname{\color{black}}%
        \expandafter\def\csname LT1\endcsname{\color{black}}%
        \expandafter\def\csname LT2\endcsname{\color{black}}%
        \expandafter\def\csname LT3\endcsname{\color{black}}%
        \expandafter\def\csname LT4\endcsname{\color{black}}%
        \expandafter\def\csname LT5\endcsname{\color{black}}%
        \expandafter\def\csname LT6\endcsname{\color{black}}%
        \expandafter\def\csname LT7\endcsname{\color{black}}%
        \expandafter\def\csname LT8\endcsname{\color{black}}%
      \fi
    \fi
      \setlength{\unitlength}{0.0500bp}%
      \ifx\gptboxheight\undefined%
        \newlength{\gptboxheight}%
        \newlength{\gptboxwidth}%
        \newsavebox{\gptboxtext}%
      \fi%
      \setlength{\fboxrule}{0.5pt}%
      \setlength{\fboxsep}{1pt}%
      \definecolor{tbcol}{rgb}{1,1,1}%
  \begin{picture}(4740.00,3600.00)%
      \gplgaddtomacro\gplbacktext{%
        \csname LTb\endcsname%%
        \put(400,343){\makebox(0,0)[r]{\strut{}\scriptsize 0}}%
        \csname LTb\endcsname%%
        \put(400,707){\makebox(0,0)[r]{\strut{}\scriptsize 10}}%
        \csname LTb\endcsname%%
        \put(400,1071){\makebox(0,0)[r]{\strut{}\scriptsize 20}}%
        \csname LTb\endcsname%%
        \put(400,1436){\makebox(0,0)[r]{\strut{}\scriptsize 30}}%
        \csname LTb\endcsname%%
        \put(400,1800){\makebox(0,0)[r]{\strut{}\scriptsize 40}}%
        \csname LTb\endcsname%%
        \put(400,2165){\makebox(0,0)[r]{\strut{}\scriptsize 50}}%
        \csname LTb\endcsname%%
        \put(400,2529){\makebox(0,0)[r]{\strut{}\scriptsize 60}}%
        \csname LTb\endcsname%%
        \put(400,2893){\makebox(0,0)[r]{\strut{}\scriptsize 70}}%
        \csname LTb\endcsname%%
        \put(400,3258){\makebox(0,0)[r]{\strut{}\scriptsize 80}}%
        \csname LTb\endcsname%%
        \put(454,235){\makebox(0,0){\strut{}\scriptsize 0}}%
        \csname LTb\endcsname%%
        \put(864,235){\makebox(0,0){\strut{}\scriptsize 100}}%
        \csname LTb\endcsname%%
        \put(1275,235){\makebox(0,0){\strut{}\scriptsize 200}}%
        \csname LTb\endcsname%%
        \put(1685,235){\makebox(0,0){\strut{}\scriptsize 300}}%
        \csname LTb\endcsname%%
        \put(2096,235){\makebox(0,0){\strut{}\scriptsize 400}}%
        \csname LTb\endcsname%%
        \put(2507,235){\makebox(0,0){\strut{}\scriptsize 500}}%
        \csname LTb\endcsname%%
        \put(2917,235){\makebox(0,0){\strut{}\scriptsize 600}}%
        \csname LTb\endcsname%%
        \put(3328,235){\makebox(0,0){\strut{}\scriptsize 700}}%
        \csname LTb\endcsname%%
        \put(3738,235){\makebox(0,0){\strut{}\scriptsize 800}}%
        \csname LTb\endcsname%%
        \put(4149,235){\makebox(0,0){\strut{}\scriptsize 900}}%
        \csname LTb\endcsname%%
        \put(4559,235){\makebox(0,0){\strut{}\scriptsize 1000}}%
      }%
      \gplgaddtomacro\gplfronttext{%
        \csname LTb\endcsname%%
        \put(4196,3153){\makebox(0,0)[r]{\strut{}\scriptsize Random Walk}}%
        \csname LTb\endcsname%%
        \put(4196,3028){\makebox(0,0)[r]{\strut{}\scriptsize PIG}}%
        \csname LTb\endcsname%%
        \put(4196,2903){\makebox(0,0)[r]{\strut{}\scriptsize EPIG}}%
        \csname LTb\endcsname%%
        \put(4196,2778){\makebox(0,0)[r]{\strut{}\scriptsize TPIG}}%
        \csname LTb\endcsname%%
        \put(4196,2653){\makebox(0,0)[r]{\strut{}\scriptsize ETPIG}}%
        \csname LTb\endcsname%%
        \put(93,1800){\rotatebox{-270.00}{\makebox(0,0){\strut{}\footnotesize KL-Divergence}}}%
        \csname LTb\endcsname%%
        \put(2507,75){\makebox(0,0){\strut{}\footnotesize TimeStep}}%
        \csname LTb\endcsname%%
        \put(2507,3419){\makebox(0,0){\strut{}\small Model Accuracy: No-TV Simulation}}%
      }%
      \gplbacktext
      \put(0,0){\includegraphics[width={237.00bp},height={180.00bp}]{gnuplot_figs/Model_Accuracy_No-TV}}%
      \gplfronttext
    \end{picture}%
  \endgroup
  \caption{Model accuracy (KL-divergence): No-TV simulation.}
  \label{Fig:MA0TV}
\end{figure}

\subsection{Effect of Epsilon}

We tested $\epsilon$ values in increments of 20\% for both EPIG and
ETPIG.  In the 4-TV simulation, higher $\epsilon$ reduces
distractions for both algorithms.  For ETPIG, all $\epsilon$ values
converge to similar learning rates by time step~250.  In the 1-TV and
No-TV simulations, lower $\epsilon$ values yield better long-term
performance for ETPIG, suggesting that TPIG's learned avoidance makes
heavy randomization counterproductive.  Pure TPIG ($\epsilon = 0$)
exhibits early artifacts smoothed by even small $\epsilon$ values,
suggesting that modest randomization stabilizes TPIG's learning
dynamics.  For PIG-based methods, higher $\epsilon$ is consistently
helpful since PIG lacks any mechanism to learn from its mistakes
and requires external randomization to escape distraction loops.
This asymmetry further confirms that the predict-then-verify mechanism
provides its own form of learned exploration control.

\section{Discussion}

By conditioning rewards on the measured decrease in model uncertainty,
the predict-then-verify approach transforms exploration into a
closed-loop feedback system.  The Noisy TV generates large predictions
but negligible actual learning, so Q-values for TV-related actions
naturally decay.  Because the reward converges to zero as the model
approaches the true distribution, TPIG inherits TD convergence
guarantees~\cite{Sutton1998}, unlike open-loop curiosity methods that
can be permanently misled by stochastic stimuli.

\subsection{TPIG's Learning Dynamics}

TPIG exhibits a characteristic two-phase behavior.  In the first
phase, TPIG explores broadly and somewhat inefficiently while
accumulating the transition data needed to compute meaningful rewards.
During this phase, PIG may outperform TPIG on learning rate because
PIG immediately pursues the highest-predicted-gain actions.  In the
second phase, TPIG's Q-values have stabilized enough to reliably
distinguish informative actions from noisy ones, and TPIG surpasses
PIG on all metrics.  This pattern is consistent across all 100 trials,
as the narrow error bands in the figures confirm.

The duration of this initial learning phase scales with environment
size.  In the $4 \times 4$ 4-TV environment, the crossover occurs
around time step~200; in the $5 \times 5$ 1-TV environment, around
time step~300; and in the $6 \times 6$ No-TV environment, around time
step~400.  This approximately linear scaling suggests that TPIG's startup cost
is predictable and proportional to the number of state-action pairs,
representing 20--40\% of the episode length across the environments
tested.
The learning rate $\alpha = 0.9$ and
discount factor $\gamma = 0.9$ were chosen to favor rapid adaptation;
preliminary experiments suggest the results are robust for all
$\gamma \geq 0.5$, with low $\alpha$ providing a conservative mode
that prioritizes distraction avoidance over exploration speed.

\subsection{The Role of Epsilon}

The interaction between $\epsilon$-greedy exploration and TPIG
reveals an important property.  In environments with many distractors
(4-TV), higher values of $\epsilon$ help both EPIG and ETPIG escape
distraction loops.  However, in environments with few or no
distractors (1-TV, No-TV), lower values of $\epsilon$ yield better
long-term performance for ETPIG.  This asymmetry arises because TPIG
already learns to avoid distractors on its own; adding random
exploration on top of an already-effective avoidance strategy simply
introduces noise that slows convergence.

For PIG-based methods (EPIG), higher $\epsilon$ is consistently
helpful in environments with distractors, since PIG lacks any
mechanism to learn from its mistakes and requires external
randomization to escape.

\subsection{Comparison with Other Approaches}

Unlike RND~\cite{Burda2019RND} or ICM~\cite{Pathak2017ICM}, TPIG
requires no auxiliary networks, no explicit distractor detection, and
no prior knowledge of which states are noisy.  Its noise resistance
emerges naturally from the closed-loop reward structure and also
improves exploration in clean environments, as the No-TV results
demonstrate.  Because the predict-then-verify mechanism is independent
of the specific divergence measure, practitioners could substitute
MSE or other metrics for KL-divergence with comparable results.

\subsection{Limitations}

The current study has several limitations.  First, the grid-world
environments are tabular.  Scaling to continuous state spaces would
require function approximation for the transition model and Q-values.
Second, our unbounded noise model (infinitely many unique TV states)
represents a worst case for novelty-based methods; real-world
distractors with finite state spaces would be easier to overcome.
Third, we do not combine TPIG's intrinsic reward with extrinsic
rewards for goal-directed tasks.
However, the $\pm 1\sigma$ bands in all figures confirm low
inter-trial variance, indicating that the results reported here
are statistically robust despite these limitations.

\section{Conclusion}

We introduced the predict-then-verify principle for curiosity-driven
exploration, instantiated as TPIG.  Across three grid-world
environments with zero to four distractors, TPIG outperformed PIG on
all metrics: fewer distractions, faster exploration, and more accurate
models.  Because
the principle depends on measured improvement rather than a specific
divergence measure, it could equally be instantiated with MSE or other
metrics, and integrated into deep RL curiosity methods such as
ICM~\cite{Pathak2017ICM} or RND~\cite{Burda2019RND}.  The choice of divergence measure becomes a tunable design parameter
rather than a fundamental algorithmic decision.  Future work includes
scaling to continuous environments, testing with MSE and other
divergence measures, and combining intrinsic and extrinsic rewards.


\bibliographystyle{flairs}
\bibliography{Thesis_David_Mathews}

\end{document}
